\documentclass[]{ANStopical}

%%% Packages Required by Class (already included)
% fancyhdr
% lastpage
% titling
% titlesec
% ragged2e
% enumitem
% amsmath
% graphicx
% geometry
% newtxtext
% newtxmath
% hyperref
% cleveref
% caption
% authblk
% apptools
% appendix
% ifpdf
% epstopdf

%%% Some other useful packages
\usepackage{bm}
% \usepackage{tikz}
% \usepackage{color}
% \usepackage{subcaption}
% \usepackage{algcompatible}
% \usepackage{array}

\usepackage{nomencl} % If needed

\title{Gradient-Enhanced Gaussian Process Regression for Non-Intrusive
POD-Based Reduced-Order Models: An Adjoint-Informed 1D Diffusion Study}

%%% Authors (use arabic numbers: 1, 2, 3, etc. for affiliationNumber)
\addAuthor[andrewnj03@tamu.edu]{Andrew N. Jeffery}{1}
\addAuthor[]{Jean C. Ragusa}{1}
\addAuthor[]{Jim E. Morel}{1}

%%% Affiliations (from authblk)
\addAffiliation{1}{Department of Nuclear Engineering, Texas A\&M University, College Station, TX}

%%% Write text for abstract
\Abstract{%
\textit{Abstract:}
\begin{itemize}
    \item Expensive coupled-physics solves $\to$ cheap surrogate. This is the 1D sandbox version of that problem.
    \item Pipeline: FD diffusion solve $\to$ masked QoI $\to$ POD basis $\to$ GP surrogate for the POD coefficients.
    \item New idea: an adjoint solve gives free gradients of each POD coefficient. Train a gradient-enhanced GP (GE-GPR) on those instead of a value-only GP.
    \item Test cases: layered material properties, a smooth void/Teflon interface, maybe both combined.
    \item Punchline: GE-GPR hits the same accuracy with [???]$\times$ fewer training solves.
    \item Takeaway: adjoint gradients are a cheap way to make ROM surrogates more sample-efficient.
\end{itemize}
}

%%% List up to 5 keywords separated by a comma
\keywords{reduced-order modeling, proper orthogonal decomposition, Gaussian process regression, adjoint sensitivity, gradient-enhanced surrogate}

%%% Provide a short title for the header on odd pages
\shortTitle{Gradient-Enhanced GPR for POD-ROM --- M\&C 2027}

%%% Provide a short author listing for the header on even pages
\authorHead{A. N. Jeffery, J. C. Ragusa, J. E. Morel}

\begin{document}

\section{Introduction}\label{sec:intro}

\begin{itemize}
    \item Expensive coupled-physics solves are too slow to call over and over (design loops, UQ). Non-intrusive ROMs fix that with a cheap parameter $\to$ QoI surrogate.
    \item Usual recipe: POD on a QoI snapshot matrix, project to a few coefficients $\alpha(\mu)$, fit a GP to $\mu \mapsto \alpha(\mu)$.
    \item Problem: a plain GP only learns from function values. Every expensive solve gives one data point per mode.
    \item Adjoints fix that almost for free: one adjoint solve per POD mode gives the exact gradient w.r.t.\ every parameter at once.
    \item So: feed those gradients into the GP too $\to$ gradient-enhanced GPR (GE-GPR / GEK). Cite \cite{bouhlel_gek,han_gek,ulaganathan_gek,backhaus_gek,semler_weiser_gek}.
    \item What this paper does:
    \begin{itemize}
        \item derives the adjoint sensitivities of the POD coefficients for the diffusion solver (material layers, plus a smooth moving-interface geometry)
        \item builds GE-GPR from scratch (scikit-learn has no gradient-observation support) and checks it against standard GPR
        \item runs the test problems and how gradients impact the learning
        \item lays out a verification (identity check, FD gradient check, sweep plots) that's reusable elsewhere.
    \end{itemize}
\end{itemize}


\section{Sandbox Forward Model and Baseline ROM}\label{sec:fwd}

Background: The standard, gradient-free pipeline both test cases build on.

\begin{itemize}
    \item 1D one-group diffusion, piecewise-constant material layers, Dirichlet BCs. Cell-centered 2nd-order FD $\to$ sparse system $A(\mu)\phi(\mu)=s(\mu)$
    \item Parameters $\mu$: per-layer $D$, $\Sigma_a$, $q$
    \item QoI = a masked window of the flux, $y(\mu)=H\phi(\mu)$. Snapshot matrix $\to$ thin SVD $\to$ rank-$r$ POD basis $\Psi$ by an energy cutoff
    \item Coefficients by projection: $\alpha(\mu) = \Psi^\top y(\mu)$
    \item Baseline surrogate: one GP per mode, ARD-RBF kernel (scikit-learn)
    \item Error metrics used everywhere after this: field error $\varepsilon_y$, compared against the POD floor $\varepsilon_{\mathrm{proj}}$
\end{itemize}


\section{Adjoint Sensitivities of the POD Coefficients}\label{sec:adjoint}

The gradients that feed Section~\ref{sec:gegpr}.

\begin{itemize}
    \item Standard trick: a scalar QoI $J=g^\top\phi$ comes from one adjoint solve $A^\top\lambda=g$, via $J=\lambda^\top s$
    \item Each POD coefficient $\alpha_k$ is exactly this kind of scalar QoI ($g_k=H^\top\psi_k$) $\to$  adjoint solve per mode gives $\partial\alpha_k/\partial\mu_j$ for every $\mu_j$ at once
    \item Material-layer case: clean closed forms for $\partial\alpha_k/\partial q_m$ and $\partial\alpha_k/\partial\Sigma_{a,m}$; $\partial\alpha_k/\partial D_m$ requires more work
    \item Geometry case: a hard material boundary isn't differentiable at the boundary. Fix: swap it for a smooth sigmoid switch at the one moving interface (void/Teflon), derive the chain rule through it (two knobs: interface location, transition sharpness)
    \item Same adjoint machinery works for both cases; geometry just swaps a 0/1 layer indicator for a smooth per-cell direction vector
    \item Verification, same three checks each time: identity check ($\lambda^\top s$ vs.\ direct projection, matches to solver tolerance), finite-difference gradient check, and a 1D sweep with a tangent-line overlay
\end{itemize}


\section{Gradient-Enhanced Gaussian Process Regression}\label{sec:gegpr}

The main method contribution.

\begin{itemize}
    \item Idea: train one GP per mode on both function values and adjoint gradients, not just values. A derivative of a GP is still a GP, so this is one joint covariance model.
    \item Needs an augmented kernel: value-value, value-gradient, gradient-gradient covariances (closed form for ARD-RBF). Bigger matrix ($M(1+p_{\mathrm{act}})$ instead of $M$), fit by maximizing the augmented marginal likelihood.
    \item Built from scratch, since scikit-learn's GP can't take derivative observations.
    \item Everything else about the surrogate matches the baseline (one GP per mode, same kernel family). Isolates the effect of adding gradients.
    \item Same code reused unmodified on the geometry case, just with a different sampler/adjoint front end feeding it. Might point to this demonstrating that the method doesn't care what $\mu$ physically means.
\end{itemize}


\section{Results}\label{sec:results}

\begin{itemize}
    \item \textbf{Material-layer case:} SVD spectrum/rank choice, projection-error floor vs.\ $r$, then a learning curve comparing GE-GPR vs.\ standard GPR at matched training-set sizes. Report as "GE-GPR at $M{=}\cdot$ matches standard GPR at $M{=}\cdot$" [???].
    \item \textbf{Geometry case:} Run once with only the interface location active (close to a 1D problem with nearly all the SVD energy in mode 1), and once with both parameters active. Compare GE-GPR's benefit in each. 
    \item \textbf{Combined case [???]:} Material layers + geometry together, if there's room/time. Same comparison, higher active-parameter count.
    \item \textbf{Cost accounting:} An adjoint solve reuses the same factorized matrix as the forward solve, so it's cheap per sample; the trade-off is a bigger, pricier-to-fit kernel matrix.
\end{itemize}


\section{Discussion}\label{sec:discussion}

\begin{itemize}
    \item Gradient enhancement helps most when: parameter count is small, the forward solve is the expensive part.
    \item Helps least when: the coefficient map is already easy to learn from values alone (the simple geometry case)
    \item Known limitation: the augmented kernel grows with $M(1+p_{\mathrm{act}})$. Doesn't scale to large $p$ as-is.
\end{itemize}


\section{Conclusions and Future Work}\label{sec:conclusions}

\begin{itemize}
    \item Restate the contributions from Sec.~\ref{sec:intro}, past tense, with real numbers filled in
    \item Future work: Pushing past $p_{\mathrm{act}}\in\{1,2\}$, maybe using the adjoint gradients themselves as a sensitivity/dimension-reduction screen; running the same recipe on the real coupled-physics stack; a multi-output version instead of one GP per mode.
\end{itemize}

\section*{ACKNOWLEDGEMENTS}
I have no idea, thank you Dr. Ragusa and Dr. Morel?

% Bibliography instructions for this template:
%  All bibliography references should be included in the "main.bib" file.
%  The bibliography should be generated using the "bibtex" command.

\bibliographystyle{ANStopical}
\bibliography{main}

\appendix

\section{Augmented Kernel Derivative Identities}
Full algebra behind the closed-form $K_{fd}$, $K_{dd}$ terms in Sec.~\ref{sec:gegpr}, for anyone re-implementing \texttt{GEGaussianProcess} from scratch.

\section{Verification Figures}
Overflow plot dump if space runs out in the main body: identity checks, FD gradient checks, for both test cases.

\end{document}
